\documentclass[12pt]{article}
\usepackage[utf8]{inputenc}
\usepackage[T1]{fontenc}
\usepackage{hyperref}
\title{Mayangna Language Resources for K-12 Curriculum Development}
\author{Compiled by Claude Sonnet 5.5 (claude-sonnet-5-5) for Andr\'as Kornai}
\date{}
\begin{document}
\maketitle

\section{Introduction}
Mayangna (ISO 639-3 yan; Misumalpan) is spoken in Nicaragua and Honduras. This report lists only resources that appeared in web search results. Few Mayangna textbooks or dictionaries are visible online, and I found no current catalogue of them.

\section{Primary Educational Resources}
\begin{itemize}
\item UNESCO project ``Reinforcing the Transmission of Mayangna Knowledge and Culture in the Classroom'' (2010--2016, BOSAWAS area). The search summary says it produced classroom materials with Mayangna experts, teachers and community members, including preschool and third-grade materials. \url{https://ich.unesco.org/en/project-education/reinforcing-the-transmission-of-mayangna-knowledge-and-culture-in-the-classroom-00493}
\item ``The cultural and natural heritage as a teaching resource in the Mayangna teaching practice, panamahka'': a didactic study meant to support community teachers. \url{https://camjol.info/index.php/RUC/article/view/7764}
\end{itemize}

\section{Linguistic and Reference Materials}
\begin{itemize}
\item Hale, paper on training indigenous linguists for dictionary work, spelling standardization and a planned Mayangna grammar. \url{https://lingphil.scripts.mit.edu/papers/hale/papers/hale039.pdf}
\end{itemize}

\section{Religious and Community Resources}
No religious or community-organisation resources were found in these searches.

\section{Digital and Audio Resources}
\begin{itemize}
\item ``Strengthening the Miskitu and Mayangna Languages in Nicaragua Through Digital Media'' (Global Voices, 2016). It describes a small bilingual dictionary and videos made at URACCAN. It also reports that Mayangna elders objected to publishing Mayangna-language material online, so community consent must be settled first. \url{https://rising.globalvoices.org/blog/2016/05/22/strengthening-the-miskitu-and-mayangna-languages-in-nicaragua-through-digital-media/}
\end{itemize}

\section{Curriculum Development Context}
\begin{itemize}
\item Study of linguistic attitudes toward Mayangna among scholarship students from the Nicaraguan Caribbean. \url{https://revistas.unan.edu.ni/index.php/RLL/article/view/5393}
\end{itemize}
Little published material exists online. Any curriculum work would need local partners (the Ministry of Education, URACCAN, Mayangna communities) and agreed rules on publishing.

\end{document}